\documentclass[10pt]{article}
\usepackage[T1]{fontenc}
\usepackage[utf8]{inputenc}
\usepackage{lmodern,amsmath,amssymb,booktabs,longtable,array}
\usepackage[margin=22mm]{geometry}
\usepackage{xurl}
\usepackage[colorlinks=true,urlcolor=blue]{hyperref}
\usepackage{microtype}
\hypersetup{pdftitle={Supplementary material: proof objects and reproducibility},pdfauthor={Qianli Ma, Chao Wu}}
\title{Supplementary material\\[3pt]\large The exact maximum growth factor for complete pivoting\\on real \(5\times5\) matrices}
\author{Qianli Ma\textsuperscript{1,2}, Chao Wu\textsuperscript{1}\\[3pt]
\small \textsuperscript{1}Zhejiang University\quad\textsuperscript{2}WUJIE AI\\[3pt]
\small \href{mailto:qianli.ma@zju.edu.cn}{\texttt{qianli.ma@zju.edu.cn}}\quad
\href{mailto:chao.wu@zju.edu.cn}{\texttt{chao.wu@zju.edu.cn}}}
\date{11 September 2026}
\begin{document}
\maketitle
% LaTeX body fragment. Requires amsmath, booktabs, longtable, array, and xurl.
% Provenance transcription only: no certificate replay was performed to write it.
\section*{Proof objects and reproducibility index}

This supplement identifies the mathematical modules used in the article,
their proof objects, and the scopes of the recorded exact acceptances.
Historical labels beginning R, V, B, or Round identify archived sources;
they are not additional mathematical hypotheses. File and collection
names identify manifest objects. The public repository is
\url{https://github.com/hkjtsgmc79-boop/rho5-proof}; fixed release
\texttt{v1.0.1} supplies the downloadable archives and verification guide.
Relative locations refer to
the companion analytical collection or to the named extracted archive,
as specified in S8. Hashes identify bytes; mathematical validity depends
on the analytic arguments, exact acceptances, and source-composition
rules. Large certificate payloads, including transport parts for the
historical-root archive, are supplied as release attachments. The
repository documents exact sizes, checksums and computing requirements.

\section*{Reading this revision}
The accompanying article now gives a mathematical reading guide, an
explicit complete-fibre predicate, a detailed canonical-height argument,
and a quantitative two-gap transport lemma. Its appendices expand the
canonical coordinates and the transport data and trace an actual local
certificate through its source hierarchy. This index connects those
statements with their frozen mathematical inputs; it is not itself a
replacement for the large proof objects.

Release \texttt{v1.0.1} indexes the unchanged proof data frozen in \texttt{v1.0.0}.
This manuscript revision does not alter those certificates, their
thresholds, or their identities, and it does not report a new
whole-chain mathematical replay. A completed first-phase Lean project
cold rebuild is documented separately in S8. The paper distinguishes
three responsibilities:
\begin{enumerate}
\item universal analytic implications, including complete-matrix
reconstruction and feasibility of the motions;
\item exact local verification of the supplied finite instances;
\item complete coverage and agreement of all inherited source domains.
\end{enumerate}
The third responsibility includes recursive dependencies and cannot be
inferred from file counts or hashes alone. The public verification guide
specifies which component commands execute mathematics and which later
commands compose previously checked results.

\section*{S1. The constant and attainment}

\path{supplement/alpha_exact.json} contains all 62 ascending integer
coefficients of the degree-61 polynomial and rational isolating endpoints
\(L<U\), with \(U-L=10^{-200}\). The published lower bound and its
attaining equality family are due to Chen, Edelman, and Urschel.
The identity record \path{evidence/v54/rho5_v54_alpha_identity.json}
compares all coefficients and records the endpoint signs and the
one-sign-change transformation
\((1+t)^{61}P_{61}((4+5t)/(1+t))\). This identifies the same unique
root in \((4,5)\); it is not a separate global upper-bound proof.

The retained foundational reconstruction is in
\path{supplement/foundations/}. R2 supplies coefficient and resultant
reconstruction; R5 supplies rational Krawczyk isolation; R7 proves that
the reconstructed matrix is actually completely pivoted and attains
\(\alpha\). The readable attainment record is also copied to
\path{supplement/R7_candidate_primal_cp.md}.
Within the foundational directory the corresponding files are:
\path{outputs/P5_coefficients_ascending.csv},
\path{outputs/P5_exact_verification.json},
\path{outputs/R2_free_CAS_reproduction_report.md},
\path{outputs/R5_candidate_interval_newton.md}, and
\path{outputs/R7_candidate_primal_cp.md}.
The actual entry points, preserving their original relative inputs, are
\path{work/verify_p5_exact.py},
\path{work/candidate_interval_newton.py}, and
\path{work/candidate_primal_cp_certificate.py}.
They use \path{work/rational_interval.py} and the three polynomial
files \path{work/sympy_recompute/p1.txt},
\path{work/sympy_recompute/p2.txt}, and \path{work/sympy_recompute/p3.txt}.
The primal entry already rechecks the Krawczyk enclosure; running every
entry separately is not required for every review.

\section*{S2. Early bounds, complete lift, and saturation}

The scalar and saturation sources are
\path{supplement/R8_scalar_root_bounds.md} and
\path{supplement/R8_global_chart_coverage.md}.
The complete two-step lift is R3,
\path{supplement/foundations/outputs/R3_exact_lift_chebyshev_reduction.md}.
The early bounds apply after normalization of the largest initial
entry to one. The D/X saturation used in the article preserves a
nonzero fifth-pivot magnitude and the complete earlier pivoting bands;
zero fifth pivots are handled by the early bounds.

\begin{samepage}
The preserved historical R8 saturation text is accompanied by
\path{supplement/R8_SATURATION_ERRATA_V55.md}.
A zero fifth pivot need not mean a zero fourth-stage tail, and the two
full-matrix constructions are asserted equal only at a positive-product
junction. The immutable historical text has SHA-256
\begin{quote}
\path{0ffdc19be82c27f3e7d660bb915ce1bcc051908c89c246fbea60cfe2b4ebecf0}.
\end{quote}
\end{samepage}
The corrections, analytic proof map, and bounded checks are retained in
\path{evidence/v55/ANALYTIC_AUDIT.md},
\path{evidence/v55/PROOF_DEPENDENCY_MAP.md}, and
\path{evidence/v55/LOCAL_ACCEPTANCE_RECEIPT.json}.
Those finite checks are not a replacement for the universal written
arguments or a replay of the X and negative-D certificates.

\section*{S3. Canonical representative, reconstruction, and height}

The canonical-height dependency has three distinct components.
Their historical acceptance records must not be merged into a claim
that the final composer re-executed all three.

\textbf{Canonical representative (V24/T24).}
The original package is
\path{RHO5_V24_Exact_Verification_2026-09-06.zip}.
Its six complete exclusion trees contain 228,542 nodes and 114,274
contradiction leaves, with no open leaf. The package retains the
mathematical theorem, model, exact acceptor, and nested earlier inputs.
Its SHA-256 is
\begin{quote}
\path{2478b3d6c8d9189fd508e2af4695bb2086c6628d5049c75cef759bd3e455dfa8}.
\end{quote}
Here T24 denotes the canonical target proved in this package.
Its primary statement is \path{RHO5_V24/THEOREM.md}, Sections
7--8, and the exceptional representative domain is defined in
Section 2. The nested V20 theorem, Section 7, supplies an attained
same-head wall maximum; V18, Section 5, and V16, Section 4.2,
supply the complete reconstruction at any prescribed attainable
high target. V15, Sections 2 and 4--5, retains a common framework
for all columns. Appendix B now proves compactness of the complete
fixed-head high-value fiber and attainment of its exact-target wall
minimum in E, before any strict endpoint qualifications are imposed.
The V25 handoff restates these inputs but does not replace their
original proofs. To replay the six V24 exclusions together with
the nested earlier acceptances, run \texttt{python verify.py
--jobs 4 --full-chain} in \path{RHO5_V24/}; omitting
\texttt{--full-chain} runs only V23's algebra after the new V24
checks, and does not re-execute the V22/V20 exclusions.

\textbf{Row-scale reconstruction (V25).}
The original nested object is
\path{RHO5_V26_Exact_Verification_REBUILT/V25_DEPENDENCY.zip}
inside
\path{RHO5_V26_Canonical_Global_Alpha_Exact_Verification_2026-09-07.zip}.
The SHA-256 of the nested ZIP bytes is
\begin{quote}
\path{2ee17e8eb8ff5d21716675ce86eea2cc791dc2a9454d820c90bd481e164a47ca}.
\end{quote}
The universal equivalence is analytic. Its 28 symbolic checks and
rational regressions test implementation and reconstruction; they are
not a sweep of all parameter values.

\textbf{Canonical exact height (V28).}
The preserved package is
\path{rho5_v28_exact.zip},
with SHA-256
\begin{quote}
\path{bec88339a158920d9dacb15341b02f359af83c4bf0e977721d9bdb2764fa0668}.
\end{quote}
All thirteen localization fronts were continued. The proof has
152,467 records and 208,274 contraction certificates; its three
terminal boxes in the peak neighborhood are paid by the analytic and
algebraic height module. The independent replay used the earlier
V24/V25 mathematical inputs without rerunning every older tree.
V27 entrance dictionaries and other earlier dependencies are retained
where the frozen component manifests require them.

\textbf{Additional upstream replay, 12 September 2026.}
The unmodified V24 entry with \texttt{--full-chain} was executed in
a fresh directory. All 22 V20/V22/V24 exclusion trees passed:
1,275,818 nodes, 637,920 contradiction leaves, and no unresolved
leaf, together with the invoked nested algebra and reconstruction
checks. A separate fresh V28 run passed all eleven required modules,
including all thirteen localization grafts and the remaining analytic
and algebraic interfaces. The respective elapsed times were
862.85 and 44.94 seconds; the frozen static inputs were unchanged.
Logs, identities, environment details, and precise scopes are retained
in \path{data/replay_receipts/}. These are fresh executions of the
frozen acceptors, not independently rewritten acceptors or a new
full-theorem replay; the universal analytical source-to-model
arguments remain separate mathematical responsibilities.

\textbf{Readable height argument.}
Appendix~B of the revised article expands the mathematical chain:
qualified representative existence, necessary and sufficient row-scale
reconstruction, complete localization, a constrained maximum, forced
contact equalities, and the exact algebraic height. The equality pattern
is a conclusion of the maximum argument, not an assumption on all
physical sources. A common polynomial root or an isolated stationary
point does not replace these steps.

The statements inherited by the representative theorem retain their
full head budgets, positivity qualifications, and reconstruction ranges.
The frozen older inputs supply those qualifications; the low-\(r\)
certificate consumes the resulting bound through a separately qualified
local transport. The dependency direction is canonical height to complete
\(X\), and then complete \(X\) to negative-diagonal safety. Final
negative-diagonal coverage is not used to establish its own upstream
canonical-height input.

\textbf{Canonical-height replay interfaces.}
The V28 entry \path{verify_v28_all.py} checks the hash manifest,
integer interval routines, rejection controls, and all required modules.
The component \path{verify_v28_coverage.py} attaches thirteen
\path{certs/front_*.jsonl.gz} payloads to the preserved partial
localization tree. It verifies original and child boxes, contractions,
conditional rows, and the three actual containments in the peak box.
The resulting 152,467-record cover must not be confused with its
incomplete earlier checkpoint.

The entry \path{verify_v28_symbolic_bridge.py} checks the 37-expression
guard dictionary, while \path{verify_v28_bridges.py} checks 42 rational
transport and qualification statements and their match with the
Jacobian box. Under \path{RHO5_V26_REBUILT/},
\path{verify_upper_tree_recovered.py} accepts the 169-node upper-target
tree, and \path{verify_jacobian_recovered.py} accepts the uniform
nine-variable certificate. The recovered certificates are equivalent
repaired proof objects, not a claim to reproduce lost earlier bytes.
The entries \path{verify_v28_carrier.py} and the rebuilt
\path{verify_resultant.py} check the carrier identities and resultant
against the complete defining polynomial. Root uniqueness in the
article also follows from the separately established root in $(4,5)$.

\section*{S4. The complete cross-saturated domain}

The complete X certificate collection comprises the following archive
objects. The table gives their exact filenames and replay entry points;
commands are run from the named extracted package root. Individual
archive instructions specify their input dependencies.
\begin{longtable}{@{}>{\raggedright\arraybackslash}p{.20\linewidth}>{\raggedright\arraybackslash}p{.48\linewidth}>{\raggedright\arraybackslash}p{.25\linewidth}@{}}
\toprule
Domain and source&Archive filename&Exact entry\\
\midrule
\endhead
\(0<k\le2\), V31&
\path{rho5_v31_exact.zip}&
\texttt{python} \path{verify_all.py}\\[2mm]
\(k\ge2.2\), Round43&
\path{ROUND43_MINIMAL_REPLAY.zip}&
\texttt{bash} \path{repro/replay.sh}, in \path{RHO5_CQG_ROUND43}\\[2mm]
\(2.15\le k\le2.2\), and strict high \(r\) for \(2.1\le k\le2.2\), Round44&
\path{ROUND44_MINIMAL_REPLAY.zip}&
\texttt{python3} \path{repro/replay_new.py} \texttt{.}\\[2mm]
\(2.14\le k\le2.15\), \(r\le k\), V33&
\path{rho5_v33_exact.zip}&
\texttt{python} \path{verify_all.py}\\[2mm]
\(2.1\le k\le2.14\), \(r\le k\), Round45&
\path{RHO5_ROUND45_EXACT.zip}&
\path{rho5_round45_exact/replay.py}\\[2mm]
\(2<k\le2.1\), \(r>k\), corrected V34/Round46&
\path{RHO5_V34_HIGH_EXACT_FULL.zip}&
\path{rho5_v34_high/replay_high.py}\\[2mm]
\(\gamma/2\le k\le2.1\), \(r\le k\), \(F\ge\gamma\), Round47&
\path{RHO5_ROUND47_ALPHA_FULL.zip}&
\path{rho5_round47_alpha/replay.py}\\
\bottomrule
\end{longtable}
The last three ZIPs have respectively 222,151,979, 131,835,895, and
367,175,732 bytes. The last package, comprising complete roots I and II,
has SHA-256
\begin{quote}
\path{ddab67679b4faf3338beddaaea25dce99edfabeab937a0da07aa534c7a866de8}.
\end{quote}
Round46 uses the corrected exact factor-graph implementation, not the
superseded one. V36 assembles the unrestricted complete X theorem:
\path{supplement/V36_THEOREMS.md}, with its full source at
\path{rho5_v36_exact/THEOREMS.md}.
The package instructions specify Python/SymPy, C++17, and Boost
requirements. Their previously recorded acceptances are not represented
as new runs performed in preparing this index.

\textbf{Two distinct analytic entry rules.}
The low-r package's \path{source/THEORY.md}, Theorem V35-local,
and \path{source/FLOW_TUBE_THEOREM.md} describe the 22-coordinate
X-to-room wall flow. It decreases \(r+w\), then applies an actual
signed tail permutation at \(w=-r\); all sixteen strict room
conditions are certified at the endpoint. Its upstream room bound
uses the original two-face normalization (V10, Sections 2--6),
the tail bound (V11, T11), the middle-face reduction (V12), and
the independent canonical-height result in S3. The separate A4
exit invokes V27, Theorem 4.1, with all six weak cofactors.
The 24-coordinate two-gap transport in V36, Section 6, instead
preserves \(r+w\) while eliminating the two cross gaps and consumes
the complete X bound. Neither it nor the later V43 guard flow
is an upstream assumption of the low-r certificate.

\section*{S5. Six-variable maximum in the negative diagonal domain}

V36 gives the attained actual-tail maximum and qualified double-gap
flow. V37 eliminates \((\beta,p,e,r,s,t)\) at a fixed seventeen-coordinate
framework, on each nonempty high-value fibre in the normalized
\(p\ge1\) range. These are same-source feasibility statements, without
packet-rank assumptions. The preserved texts are
\path{supplement/V36_THEOREMS.md} and
\path{supplement/V37_SIX_VARIABLE_ELIMINATION.md}; the full V37 source is
\path{rho5_v37_exact/theory/SIX_VARIABLE_ELIMINATION.md}.
The complete framework model is \path{supplement/B17_FULL.json}, SHA-256
\begin{quote}
\path{39a65f2cbd6089dab3a403c6a72f365a02003737a641be86c3cf73cb70e76de3}.
\end{quote}
The model contains canonical images of every hypothetical source above
\(\alpha\), including the proper negative domain with \(k\le2\).
It is not restricted globally to a selected contact equation.

\section*{S6. Acceptance rules and their thresholds}

Let \(\gamma=4132517/10^6<\alpha\). Contradiction rules exclude complete
sources with \(F\ge\gamma\); height and qualified transport rules may
instead establish only \(F\le\alpha\). An alpha-safe source need not
be infeasible at the lower threshold. Exact nonnegative-weight,
interval, contraction, and ancestor-image checks are part of acceptance,
not merely discovery output or hash comparison.

The final local record has 47 old contradiction leaves, three alpha-safe
leaves (H, A, B14), one V53 contradiction, one B15 canonical-source
threshold exclusion, and two B16 canonical-source threshold exclusions.
Each B15/B16 source covers both qualified conditional charts. The joint
conclusion is alpha safety, not whole-anchor gamma infeasibility.

Within the complete B preservation archive, the actual mathematical
interfaces include \path{deep_math.py},
\path{alpha_cover.py}, and the frozen depth500 verifier; source assembly
uses \path{alpha_parent_cover.py} and
\path{final_overlay_accept.py}. The final local mathematical producer
\path{verify_joint_anchor.py} has SHA-256
\begin{quote}
\path{548de0119a5d69ff543ab4da4c5c7c5fd6a63e7f5707c4c5c70a7af4ad548bba}.
\end{quote}
The source composer \path{final_overlay_accept.py} has SHA-256
\begin{quote}
\path{974b79b97732b02008207fec847a8a68796fcb939c5acf862340cb6c5acc0a13}.
\end{quote}
The joint producer executes its frozen child verifiers; the later
composer consumes identified earlier mathematical receipts and checks
their source bindings. It does not re-execute all inherited mathematics.

\textbf{Quantitative two-gap transport.}
The article's expanded transport lemma and Appendix~C use the V36
\path{two_gap_certificate.json},
\path{verify_two_gap_certificate.py}, \path{gap_model.py}, and
\path{support/x_local_model.py} from \path{rho5_v36_exact.zip}.
They use twenty-three differentiation variables at fixed \(p\), with
bounds proved on the full twenty-four-dimensional box, including
variation of \(p\). All four boxes have the common parameter range
\[
\frac{726212331}{500000000}\le p\le
\frac{727012331}{500000000}.
\]
The strict entrance condition applies to the whole accepted source
region. An equality case of that entrance test is not accepted by this
lemma unless an additional proved rule handles it.

\textbf{Actual necessary-row example.}
The revised article displays a genuine 49-row contradiction on a
67-coordinate lift. The complete rows and bounds are in
\path{rho5_v53_exact/proof/TERMINAL_DETAILS.json}, inside the public
\path{rho5-final-anchor-example.zip}. Its source binding is recorded in
\path{SOURCE_CERTIFICATE.json} and
\path{DEPTH19_COMPLETE_SOURCE_PROJECTION.json} in the same directory.
Byte-identical copies of these three small inputs accompany this
revision under \path{data/certificate_example/}; the corresponding
V36 transport input is \path{data/two_gap_certificate.json}.
These extracts make the displayed arithmetic and coordinate tables
easy to inspect. Full source verification still uses the frozen
archive and its acceptors.

\section*{S7. Complete negative-diagonal coverage and binding}

The preserved base tree is
\path{rho5_independent_acceptance_20260911_11/checkpoint/B17_ROOT.json}
inside the historical-root archive. It has 2,102,766,386 bytes and SHA-256
\begin{quote}
\path{8e138dce9a65e49252a1f4714229b1ca95227943913f1eeefdf51714810a563c}.
\end{quote}
Its exact cold replay accepted 749,693 nodes: 374,846 splits,
374,839 contradiction terminals, four alpha-safe terminals, and four
explicit old-rule open terminals. The 750 enclosed B44 continuations
have no internal open terminal. The cold record
\path{rho5_independent_acceptance_20260911_11/COLD_REPLAY.json}
has SHA-256
\begin{quote}
\path{55f3a9dd5090cfd9e21c479444d5cca1fc15fde2b5c22b8180bbd1c834156b37}.
\end{quote}
The recorded Fraction replay used eight workers and took 22,815.35
seconds. That is a historical mathematical acceptance, not the runtime
of the final composition.

The article's \(Q_1,Q_2,Q_3,Q_4\) are the following actual old-root
boxes; the historical indices are retained here to locate their payloads:
\begin{center}
\begin{tabular}{@{}lrll@{}}
\toprule
Article&Index&Original-root path&Paid responsibilities\\
\midrule
\(Q_1\)&338726&\texttt{10011000010101010}&111/111\\
\(Q_2\)&435003&\texttt{10111000001100011000010111011}&240/240\\
\(Q_3\)&563285&\texttt{1101000101111}&159/159\\
\(Q_4\)&675104&\texttt{1110101001111}&56/56\\
\bottomrule
\end{tabular}
\end{center}
The sum \(111+240+159+56=566\) counts constituent sources inside
these four original-parent obligations. It is distinct from the
\(465+4=469\) count of original parents. Appendix~D explains the source
hierarchy. The revision's
\path{SOURCE_ACCOUNTING_566_TO_469.json} gives a record-by-record index
into the four frozen parent receipts, together with the actual root
bindings and original file identities. It is an index of existing
records, not an additional mathematical acceptance. Relative paths
inside a transformed parent must be interpreted through their full
inherited source maps, rather than concatenated to root paths.

The last local subtree, at relative path \texttt{10011101}, covers
54 terminal positions and pays only the final one of the 240
responsibilities of parent 435003. The original-parent accounting is
465 old accepted parents plus these four completions, hence 469/469
and effective open count zero. The old root bytes still contain four
O records; they are paid by the separate cover, not relabelled.
The article combines the historical types from S6 into 48 contradiction
terminals, three height or transport terminals, and three canonical-source
threshold exclusions. These are the same 54 terminal responsibilities,
not additional proof objects.

\textbf{Final local proof.}
\path{evidence/joint/ANCHOR_COMPLETE.json} and
\path{evidence/joint/COMPLETED_ANCHOR_LEDGER.json} record the B15/B16/V53
joint acceptance. The former, from the recorded 184.37-second run,
has SHA-256
\begin{quote}
\path{648aa5322cfd518d1eaa549bf806b3daab67123078c2aa64909ba97831dd11d1}.
\end{quote}
\textbf{Last parent.}
\path{evidence/final_overlay/PARENT_435003_COMPLETE.json} has SHA-256
\begin{quote}
\path{6451a87db7383915b58bb3174581a4c81c38eca9a8aae32641d00083ad6cacb9}.
\end{quote}
\textbf{Actual root cover.}
\path{evidence/final_overlay/ROOT_COMPLETE_OVERLAY.json} has SHA-256
\begin{quote}
\path{43dbb784c871ecad4b6ff10047f39a2dad7a69a0b6eb392b40c9a56306e6ea34}.
\end{quote}
Its 357 direct dependencies are enumerated in
\path{FINAL_OVERLAY_DIRECT_DEPENDENCIES.json}. Recursive runtime and
upstream theorem dependencies are additional objects; this list alone
is not the complete proof archive.

\section*{S8. Artifact collections and reproducibility scope}

The manifest groups the proof objects into the following logical
collections. Each collection groups the retained files by its
mathematical role; the public asset names are given below.
\begin{description}
\item[Companion analytical collection.]
The relative \path{supplement/} and \path{evidence/} locations used in
S1--S7 contain the analytical sources, selected exact inputs, and
acceptance records. The foundational \path{work/} and \path{outputs/}
inputs retain their relative structure under
\path{supplement/foundations/}; their identities are listed in
\path{FOUNDATIONAL_MINIMUM_MANIFEST.json}.
\item[Canonical and complete X certificates.]
The archive objects named in S3--S4 include their frozen mathematical
inputs and acceptors. Their manifests determine the nested dependencies.
\item[Complete negative-diagonal certificates.]
The historical-root and final-cover archives below contain the base
root, its exact runtime, and the four complete source compositions.
\end{description}

The two complete negative-diagonal archives are:
\begin{description}
\item[\path{historical_root.tar.gz}]
2,189,588,699 bytes; SHA-256
\path{15ef795e85d0d295e92fdd629f313b9c284cc6519388011862aee3c1e18e3f73}.
\item[\path{final_cover.tar.gz}]
1,385,430,364 bytes; SHA-256
\path{1c9fb99d184eaf9d79b44d33826e4db8fbced8db41ee2b46ec2e1616bd12361c}.
\end{description}
Their retained copies have passed whole-archive hash checks, recorded
in the preservation collection's \path{LOCAL_ARCHIVE_RECEIPT.json}.
This verifies byte identity, not a further mathematical replay.
The second archive retains the relative locations
\path{rho5_deep500_all_open_20260910_08/},
\path{rho5_b16_v53_closure_20260911_18/},
\path{rho5_b16_20260911_17/}, and
\path{rho5_cqg_continuous_20260910/run_01/runtime/models/B17_FULL.json}.
The first contains the base tree and its
\path{runtime/R54_SOURCE_MANIFEST.json}, listing 959 runtime sources,
under \path{rho5_independent_acceptance_20260911_11/}.
These are archive-relative locations. Frozen absolute-location bindings
require an explicit replay mapping that retains every mathematical
input identity and source check.

\textbf{Final local source.}
The 18,292,611-byte final-anchor review bundle, SHA-256
\path{88a05ffe78a54d0c552aa6c5f2294e8614b7f78296efbbfb0cd5332322d21e15},
contains a standard-library-only short replay for Python 3.10 or later:
\begin{verbatim}
python -S restore_inputs.py
python -S verify_joint_anchor.py --v53-dir rho5_v53_exact \
  --b16-dir B_STRUCTURE_16_DELIVERY --output-dir replayed_anchor
\end{verbatim}
The output directory must be empty or absent. The result pays the
54-leaf anchor only; its standalone whole-parent and whole-domain
flags correctly remain false.

\textbf{Complete negative-diagonal domain.}
From the extracted historical-root directory, the exact root entry is:
\begin{verbatim}
python -B -S runtime/r54_replay.py --tree checkpoint/B17_ROOT.json \
  --out REPLAY_ROOT_NEW.json --workers 8 --backend fraction --allow-open
\end{verbatim}
The new output must not already exist. Its four expected O records
are covered by the separately accepted parent proofs.
\path{deep_math.verify_branch} and \path{alpha_cover.verify_cut} are
mathematical interfaces; \path{alpha_parent_cover.py} and
\path{final_overlay_accept.py} compose accepted sources without
rerunning all the mathematics of their input receipts. A complete
independent replay therefore includes both the parent mathematics
and the final source binding. The frozen 239/240 baseline must be
preserved when attaching the last local proof, with a hash-bound replay
configuration and fresh output locations. Python assertion optimization
must not be enabled.

\textbf{Public download and environment guide.}
The published inventory totals 4,494,382,372 bytes, including intentional
asset overlap. This is not a transitive dependency-closure claim.
\path{docs/DOWNLOADS.md} separates download size, extracted size and
working-space estimates. The helper \path{scripts/fetch.py} verifies
asset hashes and reconstructs the multipart root archive; downloading
is separate from mathematical acceptance.

C++ components require a compatible C++17 compiler and the specified
Boost headers. Python dependencies, working directories and migration
constraints are given in \path{docs/VERIFY.md} and
\path{docs/VERIFY.zh-CN.md}. The global theorem also uses S1--S5,
including the complete X certificates, attained lower bound and
analytic boundary arguments. \path{THEOREM_CHAIN.md} and
\path{REPRODUCIBILITY.md} specify dependencies, entry points and scopes.
Recorded component acceptances and source compositions are preserved;
a tested, new-machine single-command replay of the entire theorem is
not claimed. Publication and hash checks do not constitute that replay.

\textbf{First-phase Lean project cold rebuild.}
The frozen project was rebuilt on Linux on 12 September 2026, starting
with no compiled Rho5 project objects. All 634 product modules,
three additional example/audit modules, and 47 named axiom checks
passed. The toolchain was Lean 4.30.0 with Mathlib revision
\path{c5ea00351c28e24afc9f0f84379aa41082b1188f}. Pinned third-party
caches were reused and missing dependency modules were rebuilt;
this was a project cold rebuild, not a new installation and rebuild
of the entire dependency environment.

The accepted route used an explicit dependency-order driver with
per-module Lean invocations. The originally documented
\texttt{lake build +Rho5.PhaseOne:olean} command failed in the
cold state; a configuration-only repair was not counted as a completed
build. The retained records distinguish that failure from the
successful alternative route. The final \texttt{\#check}/\texttt{\#print}
output still has the type
\begin{quote}
\texttt{XGlobalSafety}\(\ \longrightarrow\ \)
\texttt{RootEndpointSafety}\(\ \longrightarrow\ \)
\texttt{rho5Trace = alpha}.
\end{quote}
The 47 axiom lists agree with the frozen records. The only axioms
appearing in them are \texttt{propext}, \path{Classical.choice}, and
\path{Quot.sound}. This check does not discharge the two explicit
safety parameters. The large tree and certificate stage was not
replayed by this build.

The frozen project source ZIP has SHA-256
\begin{quote}
\path{ddb1e4300ed03464024f2a2a227cee060b6b38de96a02043f9233f0412e28f56}.
\end{quote}
The 107,995-byte cold-build receipt archive has SHA-256
\begin{quote}
\path{c0ec1be2dc8b2aff26f987d13be34a20f0f2e92a4d2389523e9fe88c9c980415}.
\end{quote}
This revision includes that archive and the completion summary under
\path{data/lean_cold_build/}. It preserves the source/object manifest,
finalization and axiom/type logs, and failed-build evidence; it is not
a standalone replay package and does not contain all module logs,
compiled objects, dependency caches, or the successful build drivers.

\end{document}
